\begin{document}
\maketitle

\begin{abstract}
This paper isolates the primitive boundary kernel. Starting from the minimal operational seed and
the one-sided boundary datum, it reconstructs the exact double, the deck involution, the Calderon
polarization, the primitive admissibility complex, the primitive seam generator, the winding
normalization $[u_\Sigma]=1$, and the normalization $\cthree=1/(8\pi)$. No Standard-Model,
phenomenological, gravitational, cosmological, or E8 claim is made in this paper.
\end{abstract}

\begin{tfptfrontbox}
\textbf{Inputs from previous papers.} Only the orientation map of Paper 0, if read first.

\textbf{New theorem contribution.} The primitive kernel
\[
\Tker^\partial=(\mathcal A,\mathcal H,D,J,\Gamma,\taudbl,\iC,\Pprim,[u_\Sigma],\cthree)
\]
is reconstructed from the one-sided boundary datum rather than inserted later.

\textbf{Not claimed here.} No carrier $3+2$ theorem, no Standard-Model gauge group, no $\alpha$,
no flavor, no gravity, no cosmology, and no E8 scale grammar.

\textbf{Falsification or audit surface.} The paper fails if the one-sided datum does not determine
the doubled datum, the Calderon polarization, the primitive Hodge selector, or the normalization of
the primitive seam class.
\end{tfptfrontbox}


\begin{tfptclaimcontract}
\textbf{Claim.} $\mathfrak S_{\min}\Rightarrow\mathcal B_{\min}\Rightarrow\Tbound\Rightarrow(\taudbl,\iC,\Pprim,[u_\Sigma],\cthree)$.\par
\textbf{Inputs.} One-sided admissible boundary datum and stated boundary-elliptic regularity.\par
\textbf{First assumptions.} Product collar, ellipticity, self-adjoint boundary realization, Calderon uniqueness, orientation normalization.\par
\textbf{Proof status.} Core theorem target for the primitive kernel.\par
\textbf{Kill condition.} Failure of the boundary datum to determine $P_C$, $\iC$, $\taudbl$, $\Pprim$, $[u_\Sigma]$, or $\cthree$.\par
\end{tfptclaimcontract}

\begin{tfptstatusbox}
\textbf{\stag{L} (Lean verified).}\ None in this paper's
body. The analytic content of this paper is not formalised
in Lean.

\smallskip\noindent
\textbf{\stag{I} (Lean interfaced).}\ The algebraic shadow of
the Calder\'on polarisation has an explicit Lean target type
in the companion Lean note. The chain
\[
H = H_+ \oplus H_-
\;\Rightarrow\;
\pi^2 = \pi
\;\Rightarrow\;
\varepsilon^2 = \mathbf 1
\;\Rightarrow\;
P_\pm
\]
is Lean-proved \stag{L} \emph{in the companion note},
with $H = H_+ \oplus H_-$ packaged as
\texttt{TFPT.Carrier.BoundaryPolarization.BoundaryPolarization}.
The Lean-proved \emph{algebraic} equivalence
\[
\texttt{BoundaryPolarization}
\;\;\simeq\;\;
\texttt{CarrierInvolution}
\]
shows that the polarisation picture and the involution picture
are two presentations of the same finite-dimensional carrier
datum. The analytic step of this paper feeds this target type.

\smallskip\noindent
\textbf{\stag{M} (manuscript theorem --- the analytic
contribution of this paper).}\ The analytic Paper-1 theorem
\[
\Tbound
\;+\;
\text{boundary elliptic regularity}
\;\Rightarrow\;
H = H_+ \oplus H_-
\;\text{(i.e.\ inhabitant of \texttt{BoundaryPolarization})}
\]
is the manuscript-level theorem of the present paper.
Producing a \texttt{BoundaryPolarization} from the boundary
spectral pair $(\mathfrak T_\partial, \text{Calder\'on
operator})$ is the analytic obligation that Paper 1
\emph{discharges at \stag{M} rigour} for the carrier chain
to apply.
\end{tfptstatusbox}

\tableofcontents

\section{Claim Map}

The paper proves the first load-bearing arrow of the TFPT series:
\[
\mathfrak S_{\min}
\Rightarrow
\mathcal B_{\min}
\Rightarrow
\Tbound
\Rightarrow
(\taudbl,\iC,\Pprim,[u_\Sigma],\cthree).
\]
All later discrete, physical, and numerical structures are treated as unavailable here. The purpose
is to make the boundary primitive kernel auditable on its own.

\begin{tfptwarningbox}
Minimality is used only as a presentation-invariant defect filtration on essentialized admissible
bordisms. It is not a preference order over desired physics. In particular, the corner count and
the later $3+2$ carrier ranks are not free minimization coordinates in this paper.
\end{tfptwarningbox}

\section{One-Sided Boundary Datum}

The starting object is
\[
\Tbound=(\mathcal A_+,\mathcal H_+,D_+,J,\Gamma,B_\Sigma).
\]
The paper should keep only the source manuscript material needed to explain one-sided admissible
data, induced boundary spectral data, and the passage to a doubled closed datum. The proof burden is
that the primitive boundary structure is not chosen by hand.

\section{Exact Double and Deck Involution}

The exact double reconstructs
\[
\Tmin^{\mathrm{cl}}=(\mathcal A,\mathcal H,D,J,\Gamma,\taudbl,\iC)
\]
with $\taudbl$ as deck involution and $\iC$ as the involution induced by the Calderon polarization.
This section carries the analytic interface to Calderon projectors and the primitive relative
spectral dynamics from the source draft.

\section{Primitive Admissibility Complex}

The primitive selector is introduced before any color, determinant, family, or QFT sector:
\[
\Pprim=\Pi_{\ker\Delta_{\mathrm{prim}}}.
\]
The section should include the primitive admissibility differential, the primitive Hodge form, and
the selector upgrade logic only as far as it is needed to prove that a later full selector can
factor through $\Pprim$.

\section{Primitive Seam Generator}

The primitive seam generator records the two normalizations that are allowed to survive this paper:
\[
[u_\Sigma]=1,
\qquad
\cthree=\frac{1}{8\pi}.
\]
The winding class is a primitive boundary output. It is not yet a family-counting or flavor
transport input; those interpretations are postponed to Papers 2 and 3.

\section{Boundary Primitive Kernel}

\begin{theorem}[Primitive boundary kernel]
\label{thm:primitive-boundary-kernel}
Under the one-sided boundary hypotheses and the primitive admissibility assumptions stated above,
the minimal boundary datum canonically determines
\[
\Tker^\partial
=
(\mathcal A,\mathcal H,D,J,\Gamma,\taudbl,\iC,\Pprim,[u_\Sigma],\cthree).
\]
In particular, the Calderon polarization induces a direct-sum decomposition
\[
\mathcal H = \mathcal H_+ \oplus \mathcal H_-,
\qquad
\iC = \mathrm{id}_{\mathcal H_+} - \mathrm{id}_{\mathcal H_-},
\]
which (passed to the finite essential carrier $E$) inhabits the Lean-side target type
\texttt{BoundaryPolarization}.
\end{theorem}

\begin{proof}[Proof structure]
The proof follows the source draft in four steps: operational seed to boundary datum; boundary
datum to exact double; Calderon polarization to $\iC$ and $\Pprim$; primitive seam normalization to
$[u_\Sigma]=1$ and $\cthree=1/(8\pi)$. The paper must keep the representation-theoretic carrier
and all observable readouts out of the argument.
\end{proof}

\begin{remark}[Lean-side target]
\label{rem:lean-target-paper1}
The output of \Cref{thm:primitive-boundary-kernel} is exactly the typed datum consumed by the
companion Lean note. Restricted to the finite essential carrier $E \subset \mathcal H$, the
decomposition $E = E_- \oplus E_+$ is precisely an inhabitant of
\texttt{TFPT.Carrier.BoundaryPolarization.BoundaryPolarization}$\,K\,E$
($K = \mathbb Q$ or $\mathbb C$ depending on the working ring). Via the Lean-proved bridges
\texttt{toCalderonProjector}, \texttt{toCalderonCertificate}, and
\texttt{toCarrierInvolution} this immediately feeds the carrier-rigidity chain
of Paper~2 and the Lean note.

Equivalently (and used in some constructions), one may produce the carrier involution
$\epscar$ directly: an inhabitant of $\texttt{CarrierInvolution}\,(\mathrm{End}_K E)$. The Lean-proved
algebraic equivalence
\[
\texttt{BoundaryPolarization}\,K\,E
\;\simeq\;
\texttt{CarrierInvolution}\,(\mathrm{End}_K E)
\]
ensures that either presentation suffices.
\end{remark}

\section{Audit Protocol}

\begin{center}
\begin{tabularx}{\textwidth}{@{}p{0.24\textwidth}Y@{}}
\toprule
\textbf{Audit item} & \textbf{Question} \\
\midrule
One-sided datum & Is every primitive object reconstructed from $\Tbound$ or from standard boundary spectral data? \\
Calderon interface & Are regularity and projection hypotheses explicit? \\
Primitive selector & Does $\Pprim$ precede the later singlet and determinant projections? \\
Winding class & Is $[u_\Sigma]=1$ fixed before any family or flavor interpretation? \\
Normalization & Is $\cthree=1/(8\pi)$ derived without empirical input? \\
\bottomrule
\end{tabularx}
\end{center}

\clearpage
\section{Main Technical Development}

This section contains the main technical development assigned to this paper by the TFPT 4.5 clean split. Cross-paper background is referenced through dependency and scope boxes; extended backend material is kept in the Technical Companion.

\input{source_extracts/01_boundary_kernel_source.tex}

\section{Source Extraction Map}

\begin{tfptsourcebox}
Use \sourceDraft:
\begin{itemize}[leftmargin=1.5em]
    \item Sections 1 and 2 only as a shortened claim map.
    \item Sections 3.1 and 3.2 for primitive core, one-sided boundary datum, doubling, deck
    involution, Calderon projector, and primitive admissibility complex.
    \item Section 4.1 only for the primitive seam generator, $[u_\Sigma]=1$, and $\cthree$.
    \item Sections 8.1--8.4 only for primitive selector and upgrade logic.
    \item Sections 9.1--9.3 for operational seed, admissible bordisms, and boundary primitive
    kernel.
    \item Appendices B and D in shortened technical form.
\end{itemize}
\end{tfptsourcebox}

\begin{tfptwarningbox}
Every occurrence of carrier rigidity, Standard-Model representations, $\alpha$, gravity,
cosmology, CMB, or E8 should be deleted or converted into a forward reference.
\end{tfptwarningbox}


\begin{tfptexports}
Exports: $\taudbl$, $\iC$, $\Pprim$, $[u_\Sigma]=1$, $\cthree=1/(8\pi)$, and the boundary primitive kernel $\Tker^\partial$.
\end{tfptexports}

\section{Not Used Here}

Carrier rigidity, the Standard-Model packet, exact electromagnetic closure, flavor transport,
QFT reconstruction, gravity/metrology, cosmology, CMB targets, E8 grammar, horizons, and transient
channels are not used as inputs in this paper.

\clearpage
\input{source_extracts/bibliography_core_theory.tex}

\end{document}
